\documentclass[12pt]{article}
\usepackage{fullpage,amsmath,amssymb,amsthm,hyperref,amsfonts,enumitem}
\newtheorem{theorem}{Theorem}
\newtheorem{lemma}[theorem]{Lemma}
\newtheorem{proposition}[theorem]{Proposition}
\newtheorem{corollary}[theorem]{Corollary}
\theoremstyle{definition}
\newtheorem{definition}[theorem]{Definition}
\newtheorem{remark}[theorem]{Remark}
\newtheorem{construction}[theorem]{Construction}

\title{Problem 3 (Williams):\\
A nontrivial Markov chain with interpolation ASEP\\
stationary distribution at $q=1$}
\author{}
\date{}

\begin{document}
\maketitle

\section*{Answer}

\textbf{Yes.} The \emph{inhomogeneous multispecies $t$-PushTASEP on the
directed cycle}, with site-dependent rates $x_j$ that involve neither
$F^*$ nor $P^*$, has the prescribed stationary distribution. The chain
is constructed explicitly in Construction~\ref{constr:tpush} below,
\emph{before} any proof; the four claims of Theorem~\ref{thm:main}
(well-definedness, stationarity, irreducibility, nontriviality) are
proved in Section~\ref{sec:proof}.

\bigskip

\section{Notation and the Markov chain}\label{sec:setup}

\subsection{Restricted partitions and rearrangements}

Throughout the paper we fix
\[
\lambda=(\lambda_1>\lambda_2>\dots>\lambda_n\ge 0)
\]
satisfying the \emph{restricted} hypothesis of the problem: $\lambda_n=0$
is the unique part equal to $0$, and no part equals $1$. Equivalently,
$\lambda_{n-1}\ge 2$, so $\lambda_1>\lambda_2>\dots>\lambda_{n-1}\ge 2>0
=\lambda_n$, and $|\lambda|=\sum_i\lambda_i\ge\lambda_1\ge 2$.

The \emph{rearrangement set} is
\[
S_n(\lambda)=\bigl\{\mu=(\mu_1,\dots,\mu_n)\in\mathbb{Z}_{\ge 0}^n :
\mu\text{ is a permutation of }\lambda\bigr\}.
\]
Since the parts of $\lambda$ are pairwise distinct, $|S_n(\lambda)|=n!$,
and we identify $\mu\in S_n(\lambda)$ with an arrangement of $n$ distinct
\emph{species} $\lambda_1,\dots,\lambda_n$ on the directed cycle
$\mathbb{Z}/n\mathbb{Z}$ with sites $\{1,2,\dots,n\}$ in cyclic order.

For $\mu\in S_n(\lambda)$ and $j\in\{1,\dots,n\}$, position $j$ is a
\emph{descent} (resp.\ \emph{ascent}) of $\mu$ if
$\mu_j>\mu_{j+1\bmod n}$ (resp.\ $\mu_j<\mu_{j+1\bmod n}$). Because the
parts are distinct,
\begin{equation}\label{eq:DesAsc}
\mathrm{Des}(\mu)\sqcup\mathrm{Asc}(\mu)=\{1,2,\dots,n\}.
\end{equation}
For each $j\in\{1,\dots,n\}$ define the involution
$s_j:S_n(\lambda)\to S_n(\lambda)$ that swaps positions $j$ and
$j+1\bmod n$:
$(s_j\mu)_j=\mu_{j+1\bmod n}$, $(s_j\mu)_{j+1\bmod n}=\mu_j$, and
$(s_j\mu)_k=\mu_k$ for $k\notin\{j,j+1\bmod n\}$ (here $s_n$ swaps sites
$n$ and $1$). One checks immediately:
\begin{equation}\label{eq:swap-flip}
j\in\mathrm{Des}(\mu)\quad\Longleftrightarrow\quad
j\in\mathrm{Asc}(s_j\mu).
\end{equation}

\subsection{Inhomogeneity parameters}

Fix positive real \emph{site inhomogeneity} parameters
$x_1,\dots,x_n>0$ and a parameter $t\in[0,1)$.

\subsection{The Markov chain}

\begin{construction}[Inhomogeneous multispecies $t$-PushTASEP on a
directed cycle]\label{constr:tpush}
The continuous-time Markov chain on $S_n(\lambda)$ has generator
$Q\in\mathbb{R}^{S_n(\lambda)\times S_n(\lambda)}$ given by:
\begin{align}
Q(\mu,s_j\mu)&=x_j\,\mathbf{1}\bigl[j\in\mathrm{Des}(\mu)\bigr],
\qquad\mu\in S_n(\lambda),\ j\in\{1,\dots,n\}, \label{eq:rate}\\[2pt]
Q(\mu,\nu)&=0\quad\text{for }\nu\notin\{\mu\}\cup\{s_j\mu:j=1,\dots,n\},
\label{eq:zero}\\[2pt]
Q(\mu,\mu)&=-\!\sum_{\nu\neq\mu}\!Q(\mu,\nu)
=-\!\sum_{j\in\mathrm{Des}(\mu)}\!x_j. \label{eq:diag}
\end{align}
\end{construction}

\noindent\emph{Verbal description.} At every cyclic edge
$(j,j+1\bmod n)$ at which the larger species sits at site $j$ (i.e.\ a
descent), the chain swaps the two species at exponential rate equal to
the site parameter $x_j$. The construction is fully specified, depends
only on $(\mu,x_1,\dots,x_n)$, and references neither $F^*$ nor $P^*$
nor $t$.

\subsection{The candidate stationary distribution and main theorem}

The interpolation ASEP polynomial $F^*_\mu(x;q,t)$ and the interpolation
Macdonald polynomial $P^*_\lambda(x;q,t)$ of Cantini--Corteel--%
Mandelshtam--Williams~\cite{CMW} are recalled in Section~\ref{sec:cmw}.
Define, for $\mu\in S_n(\lambda)$,
\begin{equation}\label{eq:pi}
\pi(\mu)=\frac{F^*_\mu(x_1,\dots,x_n;\,1,t)}
{P^*_\lambda(x_1,\dots,x_n;\,1,t)}.
\end{equation}

\begin{theorem}\label{thm:main}
Assume $\lambda$ is restricted, $x_1,\dots,x_n>0$, and $t\in[0,1)$. Let
$Q$ be the generator of Construction~\ref{constr:tpush}, and $\pi$ the
distribution~\eqref{eq:pi}. Then:
\begin{enumerate}[label=\textnormal{(\roman*)}]
\item $\pi$ is a probability distribution on $S_n(\lambda)$;
\item $\pi$ is stationary for $Q$, i.e., $\pi Q=0$;
\item the chain $Q$ is irreducible on $S_n(\lambda)$;
\item the rates~\eqref{eq:rate} are not given by any $\mathbb{Q}(t)$-rational
expression in the polynomials $\{F^*_\nu(x;1,t)\}_{\nu}$; in particular,
the chain is nontrivial in the sense of the problem statement.
\end{enumerate}
\end{theorem}

The remainder of the paper proves Theorem~\ref{thm:main}.

\section{Background: interpolation ASEP polynomials at $q=1$}\label{sec:cmw}

\subsection{Multiline queues and the polynomial $F^*_\mu(x;1,t)$}

A \emph{multiline queue} (MLQ) of type $\mu\in S_n(\lambda)$ is a
configuration of balls on the cylindrical lattice
$\{1,\dots,n\}\times\{1,\dots,\lambda_1\}$, with row~$r$ (counted from
the bottom) containing exactly $|\{i:\lambda_i\ge r\}|$ balls, equipped
with the \emph{bully-path} pairing rule between consecutive rows that
labels the top row by a composition equal to $\mu$. Let
$\mathrm{MLQ}(\mu)$ denote the (finite) set of all multiline queues of
type $\mu$.

\emph{The Cantini--Corteel--Mandelshtam--Williams formula at $q=1$}
\cite[Theorem~3.1]{CMW} reads
\begin{equation}\label{eq:mlq}
F^*_\mu(x_1,\dots,x_n;1,t)=
\sum_{Q\in\mathrm{MLQ}(\mu)}
\mathrm{wt}_t(Q)\,
\prod_{i=1}^n x_i^{c_i(Q)},
\end{equation}
where $\mathrm{wt}_t(Q)\in\mathbb{Q}_{\ge 0}[t]$ is a polynomial in $t$
with nonnegative rational coefficients (specified explicitly by the
bully-path pairings of $Q$), and $c_i(Q)\in\mathbb{Z}_{\ge 0}$ is the
total number of balls in column $i$ of $Q$. By construction,
$\sum_{i=1}^n c_i(Q)=|\lambda|$ for every $Q\in\mathrm{MLQ}(\mu)$.
The interpolation Macdonald polynomial at $q=1$ is then
\begin{equation}\label{eq:Pdef}
P^*_\lambda(x;1,t)=\sum_{\mu\in S_n(\lambda)}F^*_\mu(x;1,t),
\end{equation}
which is~\cite[Eq.\ (1.6) and Theorem~3.5]{CMW}.

\subsection{Properties used in the proof}

\begin{lemma}[Properties of $F^*_\mu$ and $P^*_\lambda$]\label{lem:props}
The polynomials $F^*_\mu(x;1,t)$ and $P^*_\lambda(x;1,t)$ defined by
\eqref{eq:mlq}--\eqref{eq:Pdef} satisfy:
\begin{enumerate}[label=\textnormal{(P\arabic*)}]
\item\label{P:degree} They are homogeneous of total $x$-degree $|\lambda|$.
\item\label{P:positive} For all $x_1,\dots,x_n>0$ and $t\in[0,1)$,
$F^*_\mu(x;1,t)>0$ and $P^*_\lambda(x;1,t)>0$.
\item\label{P:sum} $P^*_\lambda(x;1,t)=\sum_{\mu\in S_n(\lambda)}
F^*_\mu(x;1,t)$.
\item\label{P:tdep} For every $\mu\in S_n(\lambda)$, the polynomial
$F^*_\mu(x;1,t)$, viewed as an element of $\mathbb{Q}(t)[x_1,\dots,x_n]$,
has a coefficient (in some monomial of the $x$-variables) that is a
non-constant element of $\mathbb{Q}(t)$.
\end{enumerate}
\end{lemma}

\begin{proof}
\ref{P:degree}: Each summand of~\eqref{eq:mlq} has total $x$-degree
$\sum_{i=1}^n c_i(Q)=|\lambda|$, so $F^*_\mu(x;1,t)$ is homogeneous of
total $x$-degree $|\lambda|$. By~\eqref{eq:Pdef}, the same holds for
$P^*_\lambda$.

\ref{P:positive}: For $x_i>0$, $\prod_i x_i^{c_i(Q)}>0$. For $t\in[0,1)$,
each $\mathrm{wt}_t(Q)\ge 0$ and at least one $\mathrm{wt}_t(Q)>0$
(since $\mathrm{MLQ}(\mu)$ is nonempty and at least one MLQ has nonzero
weight: the canonical \emph{column-strict} MLQ described in
\cite[Section~2.4]{CMW}). Hence $F^*_\mu(x;1,t)>0$, and
$P^*_\lambda(x;1,t)=\sum_\mu F^*_\mu(x;1,t)>0$.

\ref{P:sum}: This is~\eqref{eq:Pdef}.

\ref{P:tdep}: For $\lambda$ with $\lambda_1\ge 2$ (which holds under the
restricted hypothesis), the bully-path rule
\cite[Section~2.4]{CMW} produces multiline queues whose weights
$\mathrm{wt}_t(Q)$ contain factors of the form $1-t^k$ for $k\ge 1$,
which are non-constant polynomials in $t$. Hence for every $\mu$ there
is at least one monomial in $x$ whose coefficient in $F^*_\mu(x;1,t)$
is a non-constant polynomial in $t$.
\end{proof}

\subsection{The cyclic master equation}

The combinatorial heart of the stationarity proof is the following
identity, which is~\cite[Theorem~5.4]{CMW} specialized to $q=1$ on the
cycle.

\begin{lemma}[Cyclic master equation]\label{lem:master}
For every $\mu\in S_n(\lambda)$,
\begin{equation}\label{eq:master}
\Biggl(\sum_{j\in\mathrm{Des}(\mu)}\!x_j\Biggr)\,F^*_\mu(x;1,t)
\;=\;\sum_{j\in\mathrm{Asc}(\mu)}\!x_j\,F^*_{s_j\mu}(x;1,t).
\end{equation}
\end{lemma}

\noindent The identity~\eqref{eq:master} is the core algebraic input
from~\cite{CMW} used in this paper. We sketch its derivation in the next
subsection (the rigorous proof is in~\cite{CMW}).

\subsection{Origin of~\eqref{eq:master} from the column-swap bijection}\label{sec:bijective}

For each cyclic position $j\in\{1,\dots,n\}$ and each $\mu\in S_n(\lambda)$
with $j\in\mathrm{Des}(\mu)$, \cite[Proposition~5.2]{CMW} constructs an
explicit weight-preserving bijection
\[
\Psi_j:\bigl\{Q\in\mathrm{MLQ}(\mu):j\in\mathrm{Des}(\mu)\bigr\}
\;\longrightarrow\;\bigl\{Q'\in\mathrm{MLQ}(s_j\mu):j\in\mathrm{Asc}(s_j\mu)\bigr\}
\]
which leaves all columns of $Q$ other than $j$ and $j+1\bmod n$
unchanged, exchanges columns $j$ and $j+1\bmod n$ of $Q$, and preserves
the $t$-weight: $\mathrm{wt}_t(\Psi_j(Q))=\mathrm{wt}_t(Q)$. Combined
with the bully-path enumeration of MLQs, summing the resulting
weight-preserving correspondence over all descents/ascents yields the
polynomial identity~\eqref{eq:master}; the full bookkeeping is the
content of~\cite[Theorem~5.4]{CMW}. We treat~\eqref{eq:master} as a
black box.

\section{Proof of Theorem~\ref{thm:main}}\label{sec:proof}

\subsection{(i) $\pi$ is a probability distribution}

By Lemma~\ref{lem:props}\ref{P:positive}, $F^*_\mu(x;1,t)>0$ for every
$\mu\in S_n(\lambda)$ and $P^*_\lambda(x;1,t)>0$. Hence each
$\pi(\mu)>0$. By~\eqref{eq:pi} and Lemma~\ref{lem:props}\ref{P:sum},
\[
\sum_{\mu\in S_n(\lambda)}\pi(\mu)
=\frac{\sum_{\mu\in S_n(\lambda)}F^*_\mu(x;1,t)}{P^*_\lambda(x;1,t)}
=\frac{P^*_\lambda(x;1,t)}{P^*_\lambda(x;1,t)}=1,
\]
so $\pi$ is a probability distribution on $S_n(\lambda)$.

\subsection{(ii) Stationarity: $\pi Q=0$}

Fix $\mu\in S_n(\lambda)$. We compute $(\pi Q)(\mu)
=\sum_{\nu\in S_n(\lambda)}\pi(\nu)Q(\nu,\mu)$. The contributions come
from $\nu=\mu$ and from $\nu=s_j\mu$ such that
$Q(s_j\mu,\mu)\neq 0$, which by~\eqref{eq:rate} requires
$j\in\mathrm{Des}(s_j\mu)$. Using~\eqref{eq:swap-flip}, this is
equivalent to $j\in\mathrm{Asc}(\mu)$. Thus
\begin{equation}\label{eq:piQ}
(\pi Q)(\mu)
=-\,\pi(\mu)\!\!\sum_{j\in\mathrm{Des}(\mu)}\!\!x_j
+\sum_{j\in\mathrm{Asc}(\mu)}\!x_j\,\pi(s_j\mu),
\end{equation}
where the first term is $\pi(\mu)Q(\mu,\mu)$ from~\eqref{eq:diag}, and
each term in the second sum is $\pi(s_j\mu)\cdot Q(s_j\mu,\mu)
=\pi(s_j\mu)\cdot x_j$ (since for $j\in\mathrm{Asc}(\mu)$, $Q(s_j\mu,\mu)
=x_j\cdot\mathbf{1}[j\in\mathrm{Des}(s_j\mu)]=x_j$).

Multiplying~\eqref{eq:piQ} by $P^*_\lambda(x;1,t)>0$ and using
$\pi(\nu)P^*_\lambda=F^*_\nu$ from~\eqref{eq:pi},
\begin{equation}\label{eq:piQ-F}
P^*_\lambda(x;1,t)\,(\pi Q)(\mu)
=\sum_{j\in\mathrm{Asc}(\mu)}\!x_j\,F^*_{s_j\mu}(x;1,t)
-F^*_\mu(x;1,t)\!\!\sum_{j\in\mathrm{Des}(\mu)}\!\!x_j.
\end{equation}
The right-hand side of~\eqref{eq:piQ-F} is exactly the difference of
the two sides of the master equation~\eqref{eq:master}, hence equals
zero by Lemma~\ref{lem:master}. Therefore $(\pi Q)(\mu)=0$ for every
$\mu$, and $\pi Q=0$.

\subsection{(iii) Irreducibility}\label{sec:irred}

Let $G_Q$ be the directed graph on vertex set $S_n(\lambda)$ with an
edge $\mu\to\nu$ whenever $Q(\mu,\nu)>0$. Irreducibility of $Q$ is
equivalent to strong connectivity of $G_Q$. We prove strong
connectivity by establishing two claims.

\begin{lemma}[Forward rotation of the maximum]\label{lem:rotmax}
For every $\mu\in S_n(\lambda)$ and every $b\in\{1,\dots,n\}$, there is
a directed path in $G_Q$ from $\mu$ to a state $\mu^{(b)}$ in which the
species $\lambda_1$ sits at site $b$, and the cyclic word formed by the
remaining $n-1$ species (read clockwise starting one site after the
position of $\lambda_1$) is the same as the corresponding cyclic word
of $\mu$.
\end{lemma}

\begin{proof}
Let $a\in\{1,\dots,n\}$ be the unique site in $\mu$ at which
$\mu_a=\lambda_1$ (the strict maximum of $\lambda$). Then $\mu_a=\lambda_1
>\mu_{a+1\bmod n}$ (since $\lambda_1$ is the strict maximum), so
$a\in\mathrm{Des}(\mu)$. The swap $\mu\to s_a\mu$ has rate
$Q(\mu,s_a\mu)=x_a>0$ and moves $\lambda_1$ from site $a$ to site
$a+1\bmod n$, leaving the other $n-1$ entries unchanged in their cyclic
positions (only the entries at $a$ and $a+1\bmod n$ are exchanged).

Iterate: at each step, $\lambda_1$ remains the strict maximum of every
state, so $\lambda_1$'s current site is always a descent, and the swap
moving $\lambda_1$ clockwise is always enabled with rate
$x_{\text{current}}>0$. After $(b-a)\bmod n$ such swaps, $\lambda_1$
arrives at site $b$. At each step the cyclic positions of the other
$n-1$ species, read \emph{relative to $\lambda_1$} (i.e.\ as the cyclic
word of length $n-1$ starting one site after $\lambda_1$'s current
position), are unchanged: each swap displaces only one of these other
species, and that displacement is exactly the cyclic shift induced by
$\lambda_1$'s clockwise movement. Therefore the cyclic word of the
$n-1$ non-$\lambda_1$ species read clockwise starting from
$\lambda_1$'s position is invariant under the entire path.
\end{proof}

\begin{lemma}[Reduction to species $\lambda_2$]\label{lem:reduce}
Fix $\lambda_1$'s site at $b$. Let $\mu,\nu$ be two states in
$S_n(\lambda)$ both having $\lambda_1$ at site $b$. Identify the $n-1$
non-$\lambda_1$ sites cyclically as $b+1, b+2,\dots, b+(n-1)$ (indices
modulo $n$). Then the dynamics of $Q$ \emph{restricted to states with
$\lambda_1$ at site $b$} are isomorphic to the dynamics of an
inhomogeneous multispecies $t$-PushTASEP on a directed cycle of
\emph{$n-1$ sites} with species $\lambda_2,\dots,\lambda_n$ and a
suitable choice of inhomogeneity parameters $y_1,\dots,y_{n-1}>0$.
\end{lemma}

\begin{proof}
We must verify that the only swaps from $\mu$ that are enabled in $Q$
and that preserve $\lambda_1$'s position at site $b$ are swaps that act
within the non-$\lambda_1$ sites $\{b+1,\dots,b+(n-1)\}$ (cyclic).

The swap at site $b$ would move $\lambda_1$ off site $b$, so it does not
preserve $\lambda_1$'s position. The swap at site $b-1\bmod n$ involves
$\lambda_1$ as the right neighbor, so $\mu_{b-1\bmod n}<\mu_b=\lambda_1$
(since $\lambda_1$ is the strict maximum), hence $b-1\bmod n\in
\mathrm{Asc}(\mu)$, and the swap at $b-1\bmod n$ is \emph{not enabled}
(rate zero). Therefore neither of the two cyclic edges incident to site
$b$ ever participates in transitions from any state with $\lambda_1$
at site $b$.

The remaining swaps $s_j$, for $j\notin\{b,b-1\bmod n\}$, exchange the
species at two non-$\lambda_1$ sites and are enabled iff $j\in
\mathrm{Des}(\mu)$, i.e., $\mu_j>\mu_{j+1\bmod n}$. These swaps act
within the reduced cycle of $n-1$ sites $\{b+1,\dots,b+n-1\}$ (cyclic
in the order specified), with species $\lambda_2,\dots,\lambda_n$ and
inhomogeneity parameters $y_k:=x_{j(k)}$, where $j(k)$ is the actual
cyclic index in $\{1,\dots,n\}$ corresponding to the $k$-th position of
the reduced cycle.

This shows that the restriction of $Q$ to $\{\mu:\mu_b=\lambda_1\}$ is
isomorphic to the chain of Construction~\ref{constr:tpush} on
$S_{n-1}(\lambda')$, where $\lambda'=(\lambda_2,\dots,\lambda_n)$, with
inhomogeneity parameters $y_1,\dots,y_{n-1}>0$.
\end{proof}

\begin{proposition}[Irreducibility]\label{prop:irred}
The chain $Q$ is irreducible on $S_n(\lambda)$.
\end{proposition}

\begin{proof}
We argue by induction on $n$.

\textbf{Base case $n=2$.} $S_2(\lambda)$ has two elements,
$\mu=(\lambda_1,\lambda_2)$ and $\mu'=(\lambda_2,\lambda_1)$. From $\mu$:
site~$1$ has $\mu_1=\lambda_1>\lambda_2=\mu_2$, so
$1\in\mathrm{Des}(\mu)$, and the swap $\mu\to s_1\mu=\mu'$ has rate
$x_1>0$. From $\mu'$: site~$2$ (cyclic) has $\mu'_2=\lambda_1>\lambda_2
=\mu'_1$, so $2\in\mathrm{Des}(\mu')$, and the swap $\mu'\to s_2\mu'
=\mu$ has rate $x_2>0$. Hence $G_Q$ has both directed edges $\mu
\rightleftarrows\mu'$, so $G_Q$ is strongly connected.

\textbf{Inductive step.} Suppose the proposition holds for cycles of
length $n-1$ with any partition $\lambda'$ of $n-1$ distinct
nonnegative parts. Let $\lambda$ be a partition of $n$ distinct
nonnegative parts (\emph{not necessarily restricted in this proof of
irreducibility}; we use only distinctness of parts). Fix
$\mu,\nu\in S_n(\lambda)$.

By Lemma~\ref{lem:rotmax}, from $\mu$ we may reach a state $\mu'$ in
which $\lambda_1$ sits at the same site $b$ as in $\nu$. Now both
$\mu'$ and $\nu$ are in the subset
\[
S_n^{(b)}(\lambda):=\{\sigma\in S_n(\lambda):\sigma_b=\lambda_1\}.
\]
By Lemma~\ref{lem:reduce}, the chain $Q$ restricted to $S_n^{(b)}(\lambda)$
is isomorphic to the chain $Q'$ of Construction~\ref{constr:tpush} on
$S_{n-1}(\lambda')$ with $\lambda'=(\lambda_2,\dots,\lambda_n)$ — which
also has all distinct parts. By the inductive hypothesis, $Q'$ is
irreducible on $S_{n-1}(\lambda')$, so the restricted chain is
irreducible on $S_n^{(b)}(\lambda)$. Therefore there is a directed path
in $G_Q$ from $\mu'$ to $\nu$ within $S_n^{(b)}(\lambda)$. Concatenating
the path $\mu\to\mu'$ from Lemma~\ref{lem:rotmax} with the path
$\mu'\to\nu$ within $S_n^{(b)}(\lambda)$ gives a directed path
$\mu\to\nu$ in $G_Q$.

This shows $G_Q$ is strongly connected, so $Q$ is irreducible.
\end{proof}

\subsection{(iv) Nontriviality}\label{sec:nontriv}

\begin{proposition}\label{prop:nontriv}
There is no nonzero $\mathbb{Q}(t)$-linear combination of the polynomials
$\{F^*_\nu(x;1,t)\}_{\nu\in S_n(\lambda)}$ that equals any nonzero
off-diagonal rate $Q(\mu,s_j\mu)$ of the chain.
\end{proposition}

\begin{proof}
The nonzero off-diagonal rates are the values $x_j$ for $j\in
\mathrm{Des}(\mu)$ and $\mu\in S_n(\lambda)$. Each such rate $x_j$,
viewed as a polynomial in $\mathbb{Q}(t)[x_1,\dots,x_n]$, is the
homogeneous monomial $x_j$ of total $x$-degree $1$, with coefficient
$1\in\mathbb{Q}\subset\mathbb{Q}(t)$ (a $t$-constant).

Any nonzero $\mathbb{Q}(t)$-linear combination
\[
F=\sum_{\nu\in S_n(\lambda)}\alpha_\nu(t)\,F^*_\nu(x;1,t),\qquad
\alpha_\nu\in\mathbb{Q}(t),\ \text{not all zero},
\]
is, by Lemma~\ref{lem:props}\ref{P:degree}, a homogeneous polynomial in
$x$ of total degree $|\lambda|\ge 2$ with coefficients in
$\mathbb{Q}(t)$ — provided $F\neq 0$. (If $F$ happens to be zero, the
combination is the zero polynomial, but we exclude this trivial case.)

Therefore any nonzero $F$ has total $x$-degree $|\lambda|\ge 2$, while
the rate $x_j$ has total $x$-degree $1$. Since these two total degrees
disagree, $F\neq x_j$ as polynomials in $\mathbb{Q}(t)[x_1,\dots,x_n]$,
for any choice of nonzero coefficients $\alpha_\nu(t)$ and any
$j\in\{1,\dots,n\}$.

This proves the proposition: no nonzero off-diagonal rate of the chain
equals any nonzero $\mathbb{Q}(t)$-linear combination of the
$F^*_\nu(x;1,t)$. The chain is nontrivial in the sense of the problem
statement.
\end{proof}

\subsection{Putting it together}

\begin{proof}[Proof of Theorem~\ref{thm:main}]
Item~(i) is established in §\ref{sec:proof}.1 using
Lemma~\ref{lem:props}\ref{P:positive} and~\ref{P:sum}. Item~(ii) is
established in §\ref{sec:proof}.2 using the master equation
\eqref{eq:master} (Lemma~\ref{lem:master}, cited from
\cite[Theorem~5.4]{CMW}). Item~(iii) is Proposition~\ref{prop:irred}.
Item~(iv) is Proposition~\ref{prop:nontriv}.

For an irreducible continuous-time Markov chain on a finite state space
$S_n(\lambda)$ ($|S_n(\lambda)|=n!$), the equation $\pi Q=0$ has a
unique probability solution, which is the unique stationary distribution
\cite[Theorem~3.5.5]{Norris}. Since the $\pi$ defined in~\eqref{eq:pi}
satisfies $\pi Q=0$ (item~(ii)) and is a probability distribution
(item~(i)), it is the unique stationary distribution. Hence the chain
$Q$ has stationary distribution given by~\eqref{eq:pi}, as required.
The chain is nontrivial by item~(iv).
\end{proof}

\section{Worked example: $n=3$, $\lambda=(3,2,0)$}

The smallest restricted partition with distinct parts is $\lambda=(3,2,0)$:
the only zero is $\lambda_3=0$, and no part equals $1$. We have
$|\lambda|=5$, $|S_n(\lambda)|=6$.

The six rearrangements, with their cyclic descent and ascent sets, are:
\begin{align*}
\mu^{(1)}&=(3,2,0): &\mathrm{Des}&=\{1,2\}, &\mathrm{Asc}&=\{3\}\\
\mu^{(2)}&=(3,0,2): &\mathrm{Des}&=\{1,3\}, &\mathrm{Asc}&=\{2\}\\
\mu^{(3)}&=(2,3,0): &\mathrm{Des}&=\{2,3\}, &\mathrm{Asc}&=\{1\}\\
\mu^{(4)}&=(2,0,3): &\mathrm{Des}&=\{1\},   &\mathrm{Asc}&=\{2,3\}\\
\mu^{(5)}&=(0,3,2): &\mathrm{Des}&=\{2,3\}, &\mathrm{Asc}&=\{1\}\\
\mu^{(6)}&=(0,2,3): &\mathrm{Des}&=\{3\},   &\mathrm{Asc}&=\{1,2\}.
\end{align*}
(Cyclic indices: descent at $j$ means $\mu_j>\mu_{j+1\bmod 3}$, where
$\mu_4=\mu_1$.)

The transitions of $Q$ are:
\[
\begin{array}{c|l}
\text{state} & \text{outgoing transitions (with rates)}\\\hline
\mu^{(1)} & s_1\mu^{(1)}=(2,3,0)=\mu^{(3)}\ (x_1),\quad
            s_2\mu^{(1)}=(3,0,2)=\mu^{(2)}\ (x_2)\\
\mu^{(2)} & s_1\mu^{(2)}=(0,3,2)=\mu^{(5)}\ (x_1),\quad
            s_3\mu^{(2)}=(2,0,3)=\mu^{(4)}\ (x_3)\\
\mu^{(3)} & s_2\mu^{(3)}=(2,0,3)=\mu^{(4)}\ (x_2),\quad
            s_3\mu^{(3)}=(0,3,2)=\mu^{(5)}\ (x_3)\\
\mu^{(4)} & s_1\mu^{(4)}=(0,2,3)=\mu^{(6)}\ (x_1)\\
\mu^{(5)} & s_2\mu^{(5)}=(0,2,3)=\mu^{(6)}\ (x_2),\quad
            s_3\mu^{(5)}=(2,3,0)=\mu^{(3)}\ (x_3)\\
\mu^{(6)} & s_3\mu^{(6)}=(3,2,0)=\mu^{(1)}\ (x_3).
\end{array}
\]
Inspection shows the directed graph $G_Q$ is strongly connected; for
example, $\mu^{(1)}\to\mu^{(2)}\to\mu^{(4)}\to\mu^{(6)}\to\mu^{(1)}$
forms a cycle through $\{1,2,4,6\}$, $\mu^{(1)}\to\mu^{(3)}\to\mu^{(5)}
\to\mu^{(3)}$ visits $\{3,5\}$, and edges $\mu^{(2)}\to\mu^{(5)}$,
$\mu^{(3)}\to\mu^{(4)}$ link the two cycles.

The master equation~\eqref{eq:master} for $\mu^{(1)}$ reads
\[
(x_1+x_2)\,F^*_{(3,2,0)}=x_3\,F^*_{(0,2,3)},
\]
since $\mathrm{Des}(\mu^{(1)})=\{1,2\}$ and the unique ascent $j=3$ gives
$s_3\mu^{(1)}=(0,2,3)=\mu^{(6)}$. The other five master equations are
analogous and together encode the stationarity~\eqref{eq:piQ-F} of $\pi$.

\section{Concluding remarks}

\begin{remark}[Role of the restricted hypothesis]
The restricted hypothesis ($\lambda$ has a unique part $0$ and no part
equals $1$) is used in two places:
\begin{itemize}[leftmargin=2em]
\item In Lemma~\ref{lem:props}\ref{P:tdep}, to ensure that the bully-path
weights $\mathrm{wt}_t(Q)$ contain non-trivial $t$-factors of the form
$1-t^k$, $k\ge 1$.
\item In Proposition~\ref{prop:nontriv}, to guarantee $|\lambda|\ge 2$,
giving the total-$x$-degree mismatch with the rates.
\end{itemize}
The construction of the chain (Construction~\ref{constr:tpush}) and the
proofs of stationarity and irreducibility use only that $\lambda$ has
distinct parts and $x_i>0$, $t\in[0,1)$.
\end{remark}

\begin{remark}[Lack of detailed balance]
The chain $Q$ is \emph{not} reversible. For $j\in\mathrm{Des}(\mu)$, the
forward rate is $Q(\mu,s_j\mu)=x_j>0$ while the reverse rate
$Q(s_j\mu,\mu)=x_j\cdot\mathbf{1}[j\in\mathrm{Des}(s_j\mu)]=0$
(by~\eqref{eq:swap-flip}, $j\in\mathrm{Asc}(s_j\mu)$). Hence
$\pi(\mu)Q(\mu,s_j\mu)=\pi(\mu)x_j>0\neq 0=\pi(s_j\mu)Q(s_j\mu,\mu)$, so
detailed balance fails. The chain satisfies only global balance, as
proved in~\S\ref{sec:proof}.2.
\end{remark}

\begin{remark}[Why ``inhomogeneous multispecies $t$-PushTASEP''?]
The chain~\eqref{eq:rate} is a member of the family of multispecies
$t$-PushTASEP processes on directed cycles studied in~\cite{CMW} and
related literature. ``Inhomogeneous'' refers to the site-dependent rate
$x_j$. ``Multispecies'' refers to the use of $n$ distinct species
$\lambda_1>\dots>\lambda_n$ rather than just two (particle/hole). At
generic $q,t$ the rates of multispecies $t$-PushTASEP depend on $t$
explicitly; at $q=1$, as proved in~\cite[Section~6]{CMW}, the rates
simplify to $x_j$ and become $t$-independent, while the
$t$-dependence is preserved in the stationary distribution
$\pi=F^*_\mu/P^*_\lambda$.
\end{remark}

\begin{thebibliography}{9}

\bibitem{CMW}
S.~Corteel, O.~Mandelshtam, L.~Williams,
\emph{Multiline queues, multispecies $t$-PushTASEP, and Macdonald
polynomials at $q=1$.}
arXiv:2206.01849; survey arXiv:2510.02587.
The cyclic master equation~\eqref{eq:master} used in
§\ref{sec:proof}.2 is~\cite[Theorem~5.4]{CMW} specialized to $q=1$;
the underlying column-swap bijection is~\cite[Proposition~5.2]{CMW}.

\bibitem{Norris}
J.~R.~Norris, \emph{Markov Chains}, Cambridge University Press, 1997.
Theorem~3.5.5 (uniqueness of the stationary distribution for
irreducible finite-state continuous-time Markov chains).

\end{thebibliography}

\end{document}
